\documentclass[11pt,a4paper]{article}
\usepackage[margin=2.25cm,headheight=22pt]{geometry}
\usepackage[T1]{fontenc}
\usepackage[utf8]{inputenc}
\usepackage[spanish,es-nodecimaldot]{babel}
\usepackage{lmodern}
\usepackage{microtype}
\usepackage{xcolor}
\usepackage{array}
\usepackage{tabularx}
\usepackage{enumitem}
\usepackage{fancyhdr}
\usepackage{hyperref}
\usepackage{xurl}
\usepackage{titlesec}
\usepackage{booktabs}

\definecolor{azul}{HTML}{18324B}
\definecolor{turquesa}{HTML}{087E83}
\definecolor{gris}{HTML}{52616B}
\definecolor{fondo}{HTML}{EEF4F6}
\definecolor{linea}{HTML}{CFDCE1}
\definecolor{alerta}{HTML}{A2472E}
\definecolor{seguro}{HTML}{19724B}

\hypersetup{
  colorlinks=true,
  linkcolor=turquesa,
  urlcolor=alerta,
  pdftitle={Reporte de análisis de phishing universitario},
  pdfauthor={Juan Skobalski}
}

\pagestyle{fancy}
\fancyhf{}
\fancyhead[L]{\small\color{gris}Ciberseguridad | Unidad 4 | Preentrega 4}
\fancyhead[R]{\small\color{gris}Juan Skobalski}
\fancyfoot[C]{\small\color{gris}\thepage}
\renewcommand{\headrulewidth}{0.4pt}
\renewcommand{\headrule}{\hbox to\headwidth{\color{linea}\leaders\hrule height \headrulewidth\hfill}}

\titleformat{\section}[block]{\Large\bfseries\color{azul}}{}{0pt}{}
\titleformat{\subsection}[block]{\large\bfseries\color{turquesa}}{}{0pt}{}
\titlespacing*{\section}{0pt}{0pt}{8pt}
\titlespacing*{\subsection}{0pt}{14pt}{6pt}
\setlist[itemize]{leftmargin=1.35em,itemsep=4pt,topsep=4pt}
\setlist[enumerate]{leftmargin=1.55em,itemsep=6pt,topsep=5pt}
\setlength{\parskip}{6pt}
\setlength{\parindent}{0pt}
\renewcommand{\arraystretch}{1.28}

\newcommand{\infobox}[1]{%
  \noindent\colorbox{fondo}{\parbox{\dimexpr\textwidth-2\fboxsep}{#1}}}

\newcommand{\casebox}[2]{%
  \noindent\fcolorbox{linea}{fondo}{%
    \parbox{\dimexpr\textwidth-2\fboxsep-2\fboxrule}{%
      \textcolor{azul}{\textbf{#1}}\par\vspace{5pt}#2}}}

\newcommand{\reportpage}{\clearpage\thispagestyle{fancy}}

\begin{document}

% Carátula
\thispagestyle{empty}
\vspace*{1.0cm}
{\color{turquesa}\large C O D E R H O U S E}\par
\vspace{0.4cm}
{\color{linea}\rule{\textwidth}{1.2pt}}\par
\vspace{2.2cm}
{\color{gris}\Large CIBERSEGURIDAD}\par
\vspace{0.45cm}
{\color{azul}\Huge\bfseries Reporte de Análisis\newline de Phishing}\par
\vspace{0.45cm}
{\color{turquesa}\LARGE Caso universitario y defensa multicapa}\par
\vspace{1.0cm}
{\color{gris}\large Unidad 4 | Preentrega 4}\par

\vfill
\noindent\colorbox{fondo}{%
  \begin{minipage}{\dimexpr\textwidth-2\fboxsep}
    \vspace{0.28cm}
    \noindent\makebox[3.5cm][l]{\textcolor{azul}{\textbf{Estudiante}}}Juan Skobalski\par
    \vspace{0.18cm}\color{linea}\hrule\color{black}\vspace{0.18cm}
    \noindent\makebox[3.5cm][l]{\textcolor{azul}{\textbf{Curso}}}Ciberseguridad\par
    \vspace{0.18cm}\color{linea}\hrule\color{black}\vspace{0.18cm}
    \noindent\makebox[3.5cm][l]{\textcolor{azul}{\textbf{Contexto}}}Seguridad universitaria\par
    \vspace{0.18cm}\color{linea}\hrule\color{black}\vspace{0.18cm}
    \noindent\makebox[3.5cm][l]{\textcolor{azul}{\textbf{Fecha}}}3 de noviembre de 2026\par
    \vspace{0.28cm}
  \end{minipage}}
\vspace{0.75cm}
{\color{linea}\rule{\textwidth}{1.2pt}}\par
\vspace{0.25cm}
{\small\color{gris}Evaluación de señales de alerta, exposición, respuesta inmediata y prevención.}

\reportpage
\section*{Introducción}

Este reporte analiza un intento de phishing simulado dirigido a un profesor universitario. El mensaje utiliza un beneficio económico y un plazo limitado para inducirlo a iniciar sesión en un portal externo y entregar datos bancarios. El caso representa un ataque de \textbf{Ingeniería Social}: el atacante intenta obtener la cooperación de la víctima manipulando sus emociones y su percepción de autoridad.

\subsection*{Escenario recibido}

\casebox{Correo reportado por el profesor}{``Estimado docente, su bonificación por investigación ha sido aprobada. Inicie sesión en \nolinkurl{portal-universitario-pagos.net} para confirmar sus datos bancarios antes de las 5:00 PM para recibir el depósito hoy mismo.''}

\subsection*{Objetivo y alcance}

El análisis se limita a los cuatro resultados solicitados:

\begin{itemize}
  \item identificar al menos cuatro señales de alerta;
  \item explicar la exposición personal e institucional;
  \item establecer tres acciones inmediatas para la universidad;
  \item redactar un boletín preventivo de hasta 100 palabras.
\end{itemize}

No se abre ni se consulta el dominio sospechoso. En un incidente real se preservaría el mensaje original, incluyendo sus encabezados, para analizar el remitente, la autenticación del correo y la infraestructura utilizada.

\infobox{\textcolor{azul}{\textbf{Evaluación inicial:}} la promesa económica, el plazo urgente, el dominio externo y la solicitud de datos bancarios forman un patrón consistente con phishing orientado al robo de credenciales e información financiera.}

\subsection*{Metodología}

Se comparan los elementos del mensaje con procedimientos institucionales esperables y se aplica un enfoque de defensa multicapa: personas, correo electrónico, identidad, red, monitoreo y respuesta. Ningún control aislado se considera suficiente.

\reportpage
\section*{1. Identificación de señales de alerta}

\begin{tabularx}{\textwidth}{>{\centering\arraybackslash\bfseries\color{alerta}}p{0.7cm} >{\raggedright\arraybackslash\bfseries\color{azul}}p{3.6cm} >{\raggedright\arraybackslash}X}
  N.º & Señal & Justificación \\
  \toprule
  1 & Dominio no institucional & \nolinkurl{portal-universitario-pagos.net} no demuestra pertenencia a la universidad. Combina términos confiables para parecer un portal oficial. \\
  2 & Urgencia artificial & El límite ``antes de las 5:00 PM'' busca reducir el tiempo disponible para pensar y verificar. \\
  3 & Beneficio inesperado & La bonificación apela a la gratificación y puede hacer que la víctima ignore inconsistencias. \\
  4 & Datos sensibles & Se solicitan credenciales y datos bancarios mediante un enlace recibido por correo, procedimiento de alto riesgo. \\
  5 & Saludo genérico & ``Estimado docente'' no identifica al profesor ni aporta información específica del trámite. \\
  6 & Procedimiento atípico & Una bonificación legítima debería figurar en los sistemas administrativos oficiales y poder confirmarse sin usar el enlace recibido. \\
\end{tabularx}

\subsection*{Spoofing y apariencia confiable}

El mensaje intenta realizar \textit{spoofing} de identidad: adopta el lenguaje y la autoridad de la universidad para que la víctima lo trate como legítimo. Un logotipo, una firma o un diseño profesional no prueban el origen porque pueden copiarse. También debe diferenciarse el nombre visible de la dirección real del remitente.

\subsection*{El candado no sería suficiente}

Aunque el portal utilizara HTTPS, eso solo cifraría la conexión con ese servidor. No demostraría que el sitio pertenece a la universidad. La validación debe comenzar por el dominio institucional y por un acceso independiente al sistema oficial.

\infobox{\textcolor{alerta}{\textbf{Conclusión técnica:}} la evaluación no depende de una sola señal. La combinación de dominio externo, urgencia, recompensa y solicitud de datos sensibles permite clasificar el mensaje como un intento de phishing.}

\reportpage
\section*{2. Exposición de riesgos}

Si el profesor introduce sus datos en el portal, el incidente puede superar el fraude individual y convertirse en un acceso inicial contra la universidad.

\begin{tabularx}{\textwidth}{>{\raggedright\arraybackslash\bfseries\color{azul}}p{3.7cm} >{\raggedright\arraybackslash}X}
  Riesgo & Consecuencia posible \\
  \toprule
  Robo de credenciales & Acceso no autorizado al correo, campus virtual, almacenamiento y otros servicios institucionales. \\
  Fraude financiero & Uso indebido de datos bancarios, intentos de cobro o nuevas solicitudes fraudulentas. \\
  Secuestro de la cuenta & Cambio de contraseña, persistencia mediante sesiones activas y pérdida de control por parte del profesor. \\
  Propagación interna & Envío de nuevos correos desde una cuenta legítima para engañar a colegas, estudiantes o personal administrativo. \\
  Movimiento lateral & Intentos de acceder a sistemas adicionales, especialmente si existe reutilización de contraseñas o privilegios elevados. \\
  Filtración de información & Exposición de documentos, investigaciones, datos personales y comunicaciones internas. \\
\end{tabularx}

\subsection*{Cadena probable del incidente}

\begin{enumerate}
  \item El profesor abre el portal y entrega credenciales o datos bancarios.
  \item El atacante prueba las credenciales contra el correo y otros servicios conectados.
  \item Si logra acceder, roba información o utiliza la cuenta para enviar phishing interno más convincente.
  \item La confianza en una cuenta institucional facilita la propagación y amplía el impacto.
\end{enumerate}

\subsection*{Factores que reducen el impacto}

La autenticación multifactor (MFA), las contraseñas únicas, la revocación rápida de sesiones y el principio de menor privilegio limitan el alcance. MFA no elimina todos los ataques, pero agrega una barrera relevante si la contraseña fue robada.

\infobox{\textcolor{azul}{\textbf{Riesgo principal:}} la cuenta del profesor puede transformarse en una identidad confiable controlada por el atacante y utilizarse como punto de entrada para nuevos fraudes.}

\reportpage
\section*{3. Plan de acción inmediato}

La universidad debe actuar sobre el mensaje reportado y, al mismo tiempo, determinar si otras personas lo recibieron o interactuaron con él.

\subsection*{Paso 1. Contener y preservar}

Bloquear temporalmente el dominio, la URL y el remitente en los controles de correo, DNS o navegación; buscar copias del mensaje y ponerlas en cuarentena. Conservar el correo original con sus encabezados para mantener evidencia y extraer indicadores verificables.

\subsection*{Paso 2. Investigar y proteger identidades}

Identificar destinatarios, clics, formularios enviados e inicios de sesión anómalos. Si alguien interactuó, restablecer sus credenciales, revocar sesiones, revisar reglas de reenvío y aplicar o reforzar MFA. También deben monitorearse intentos posteriores desde la cuenta afectada.

\subsection*{Paso 3. Comunicar y prevenir}

Enviar una alerta breve al personal con el asunto, el dominio y la conducta esperada; recordar cómo reportar el mensaje y acceder a los sistemas por canales oficiales. Compartir los indicadores con soporte y seguridad para mantener el bloqueo y el seguimiento.

\subsection*{Justificación de la secuencia}

La contención reduce nuevas víctimas; la investigación determina el alcance real; y la comunicación evita que el mismo pretexto siga funcionando. Ejecutar solo una de estas acciones dejaría una parte del riesgo sin tratar.

\infobox{\textcolor{seguro}{\textbf{Resultado esperado:}} retirar el correo, impedir el acceso al portal, proteger las cuentas afectadas y advertir a la comunidad antes de que el ataque se propague.}

\reportpage
\section*{4. Defensa multicapa}

\begin{tabularx}{\textwidth}{>{\raggedright\arraybackslash\bfseries\color{azul}}p{3.2cm} >{\raggedright\arraybackslash}X}
  Capa & Medida de mitigación \\
  \toprule
  Personas & Capacitación, simulaciones autorizadas y técnica de la lupa sobre el enlace. \\
  Correo electrónico & Filtros de reputación, análisis de URL y configuración de SPF, DKIM y DMARC para reducir suplantaciones del dominio propio. \\
  Identidad & MFA, contraseñas únicas, acceso condicional y revocación de sesiones comprometidas. \\
  Red y navegación & Filtrado DNS/web para bloquear dominios confirmados como maliciosos. \\
  Monitoreo & Alertas por inicios de sesión anómalos, nuevas reglas de correo y comportamiento inusual. \\
  Respuesta & Canal visible de reporte, responsables definidos y procedimiento documentado de contención e investigación. \\
\end{tabularx}

\subsection*{Boletín preventivo para estudiantes}

\noindent\fcolorbox{turquesa}{fondo}{%
  \parbox{\dimexpr\textwidth-2\fboxsep-2\fboxrule}{%
  \textbf{La lupa sobre el enlace - 83 palabras}\par\vspace{5pt}
  Antes de abrir un enlace, coloca el cursor sobre él sin hacer clic. Compara la dirección mostrada con el dominio oficial de la universidad. En este caso, \nolinkurl{portal-universitario-pagos.net} no pertenece al dominio institucional: combina palabras confiables para parecer legítimo. Si el destino es diferente, está acortado o solicita credenciales o datos bancarios, no ingreses. Abre el portal escribiendo manualmente la dirección oficial o desde la aplicación institucional. Reporta el mensaje a Soporte. La urgencia y una promesa de dinero nunca reemplazan la verificación.}}

\subsection*{Aplicación del semáforo de riesgo}

\begin{itemize}
  \item \textcolor{seguro}{\textbf{Verde:}} comunicación esperada, dominio institucional y trámite verificable en el portal oficial.
  \item \textcolor{turquesa}{\textbf{Amarillo:}} beneficio inesperado, saludo genérico o procedimiento poco habitual; detenerse y verificar.
  \item \textcolor{alerta}{\textbf{Rojo:}} dominio externo, urgencia y solicitud de credenciales o datos bancarios; no interactuar y reportar.
\end{itemize}

\section*{Conclusión}

El correo presenta señales suficientes para clasificarlo como phishing. Su eficacia depende de combinar autoridad, recompensa y urgencia para que el profesor entregue información sin verificar el procedimiento. El riesgo incluye fraude personal, compromiso de la cuenta y propagación dentro de la universidad.

La respuesta adecuada integra personas y tecnología: contener el mensaje, investigar el alcance, proteger identidades con MFA, bloquear indicadores y comunicar el caso. La técnica más simple para el usuario sigue siendo interrumpir la acción y verificar el enlace y el trámite desde un canal institucional independiente.

\end{document}
